\begin{frame}
\titlepage
\vspace{-0.6em}
\centering
\small Exposición individual
\end{frame}

\begin{frame}{Objetivos y clasificación}
\small
\begin{columns}[T,totalwidth=\textwidth]
\column{0.48\textwidth}
\casehead{Objetivo general}

Aplicar ratios financieros a información contable para evaluar liquidez, gestión, solvencia y rentabilidad como apoyo para la toma de decisiones empresariales.

\vspace{0.35em}
\casehead{Objetivos específicos}
\begin{itemize}
  \item Identificar partidas necesarias de los estados financieros.
  \item Calcular ratios mediante casos prácticos.
  \item Interpretar cada resultado en contexto.
  \item Formular decisiones a partir del diagnóstico.
\end{itemize}

\column{0.50\textwidth}
\centering
\begin{tikzpicture}[node distance=0.36cm, every node/.style={font=\scriptsize}]
\node[draw, rounded corners=2pt, align=center, minimum width=2.45cm, minimum height=0.55cm] (a) {Estados\\financieros};
\node[draw, rounded corners=2pt, align=center, minimum width=2.10cm, minimum height=0.55cm, right=of a] (b) {Ratios\\financieros};
\node[draw, rounded corners=2pt, align=center, minimum width=2.10cm, minimum height=0.55cm, below=0.48cm of b] (c) {Interpretación};
\node[draw, rounded corners=2pt, align=center, minimum width=2.45cm, minimum height=0.55cm, left=of c] (d) {Decisión\\empresarial};
\draw[-{Stealth}] (a) -- (b);
\draw[-{Stealth}] (b) -- (c);
\draw[-{Stealth}] (c) -- (d);
\end{tikzpicture}

\vspace{0.5em}
\begin{tabular}{ll}
\toprule
Categoría & Enfoque \\
\midrule
Liquidez & Corto plazo \\
Gestión & Ciclo operativo \\
Solvencia & Financiamiento \\
Rentabilidad & Utilidad \\
\bottomrule
\end{tabular}
\end{columns}
\end{frame}
